% ch8-jordan.tex — Chapitre 8 : Forme Normale de Jordan
% Ce chapitre avancé traite de la forme canonique de Jordan :
% endomorphismes nilpotents, sous-espaces caractéristiques,
% blocs de Jordan, théorème de Jordan, exponentielle de matrice,
% et applications aux systèmes différentiels.

\chapter{Forme Normale de Jordan}\label{ch:jordan}

\section*{Introduction}
\addcontentsline{toc}{section}{Introduction}

Au 
ef{ch:valeurs-propres}, nous avons vu qu'un endomorphisme $f$ d'un
$\K$-espace vectoriel $E$ de dimension finie est diagonalisable si et
seulement si ses sous-espaces propres engendrent~$E$. Malheureusement,
de nombreuses matrices ne sont pas diagonalisables~: par exemple, la
matrice
\[
  A=\begin{pmatrix}2&1\\0&2\end{pmatrix}
\]
possède la valeur propre double $\lambda=2$, mais son sous-espace propre
$E_2(A)$ est de dimension~$1$. Il est donc impossible de trouver une base
de~$\R^2$ formée de vecteurs propres de~$A$.

\medskip

Se pose alors une question naturelle~: \emph{quelle est la forme
matricielle la plus simple que l'on puisse atteindre par changement de
base pour un endomorphisme quelconque~?} La réponse est donnée par la
\emph{forme normale de Jordan}\index{forme normale de Jordan}, qui
constitue la forme canonique la plus fine pour les endomorphismes sur un
corps algébriquement clos (typiquement~$\C$).

\medskip

La forme de Jordan\index{Jordan, Camille} généralise la diagonalisation~:
une matrice diagonalisable a une forme de Jordan diagonale, tandis qu'une
matrice non diagonalisable acquiert des~$1$ sur la sur-diagonale,
regroupés en \emph{blocs de Jordan}. Cette structure encode très
précisément le défaut de diagonalisabilité.

\medskip

Les applications sont considérables~:
\begin{itemize}
  \item \textbf{Systèmes différentiels~:}\index{systèmes différentiels}
    le calcul de $e^{tA}$ pour une matrice non diagonalisable repose sur
    la forme de Jordan.
  \item \textbf{Puissances de matrices~:}\index{puissances de matrices}
    $A^k$ se calcule explicitement via les blocs de Jordan.
  \item \textbf{Stabilité~:}\index{stabilité} la classification des
    points d'équilibre de systèmes dynamiques linéaires fait intervenir
    la forme de Jordan.
  \item \textbf{Classification des endomorphismes~:} deux matrices sont
    semblables si et seulement si elles ont la même forme de Jordan.
\end{itemize}

\medskip

Dans tout ce chapitre, sauf mention contraire, $E$ désigne un
$\C$-espace vectoriel de dimension finie $n\geq 1$, et
$f\in\End(E)$.\index{corps algébriquement clos}


% ============================================================
\section{Endomorphismes nilpotents}\label{sec:nilpotent}
% ============================================================

La théorie de Jordan repose de manière essentielle sur la compréhension
des endomorphismes nilpotents. Ce sont les endomorphismes dont une
certaine puissance s'annule.

\begin{definition}{Endomorphisme nilpotent}{def:nilpotent}
  Un endomorphisme $f\in\End(E)$ est dit
  \emph{nilpotent}\index{nilpotent!endomorphisme} s'il existe un entier
  $p\geq 1$ tel que $f^p=0$ (l'endomorphisme nul). Une matrice
  $A\in\Mat_n(\C)$ est dite \emph{nilpotente}\index{nilpotent!matrice}
  s'il existe $p\geq 1$ tel que $A^p=0$.
\end{definition}

\begin{definition}{Indice de nilpotence}{def:indice-nilpotence}
  Soit $f\in\End(E)$ un endomorphisme nilpotent. Le plus petit entier
  $\nu\geq 1$ tel que $f^\nu=0$ est appelé l'\emph{indice de
  nilpotence}\index{indice de nilpotence} de~$f$. On écrit $\nu(f)=\nu$.
\end{definition}

\begin{proposition}{Propriétés des endomorphismes nilpotents}{prop:nilpotent-prop}
  Soit $f\in\End(E)$ un endomorphisme nilpotent d'indice~$\nu$, avec
  $\dim E=n$. Alors~:
  \begin{enumerate}[label=(\roman*)]
    \item $\nu\leq n$.
    \item La seule valeur propre de $f$ est~$0$, c'est-à-dire
      $\Spec(f)=\set{0}$.\index{spectre!endomorphisme nilpotent}
    \item $\det(f)=0$ et $\tr(f)=0$.
    \item Le polynôme caractéristique de $f$ est $\chi_f(X)=X^n$.
    \item Le polynôme minimal de $f$ est $\mu_f(X)=X^\nu$.
  \end{enumerate}
\end{proposition}

\begin{proof}
  (i) Considérons la suite de sous-espaces emboîtés~:
  \[
    \set{\vzero}=\Ker(f^0)\subseteq\Ker(f)\subseteq\Ker(f^2)
    \subseteq\cdots\subseteq\Ker(f^\nu)=E.
  \]
  Chaque inclusion $\Ker(f^k)\subseteq\Ker(f^{k+1})$ est large, et
  l'inclusion est stricte pour $k=0,1,\ldots,\nu-1$ (sinon $f^k=0$
  contredirait la minimalité de~$\nu$). Donc
  $\dim\Ker(f^{k+1})\geq\dim\Ker(f^k)+1$, ce qui donne
  $n=\dim\Ker(f^\nu)\geq\nu$.

  (ii) Si $\lambda$ est valeur propre de $f$, il existe $\vec{v}\neq\vzero$
  tel que $f(\vec{v})=\lambda\vec{v}$. Alors
  $f^\nu(\vec{v})=\lambda^\nu\vec{v}=\vzero$, donc $\lambda^\nu=0$ et
  $\lambda=0$.

  (iii) Le déterminant est le produit des valeurs propres (comptées avec
  multiplicité) sur~$\C$, et la trace est leur somme. Puisque la seule
  valeur propre est~$0$, on obtient $\det(f)=0$ et $\tr(f)=0$.

  (iv) Le polynôme caractéristique est de degré~$n$ et ses racines (dans~$\C$)
  sont les valeurs propres. Puisque $0$ est la seule valeur propre, de
  multiplicité algébrique~$n$, on a $\chi_f(X)=X^n$.

  (v) Le polynôme minimal divise le polynôme caractéristique $X^n$,
  donc $\mu_f(X)=X^m$ pour un certain $m\leq n$. Par définition,
  $f^m=0$ et $f^{m-1}\neq 0$, donc $m=\nu$.
\end{proof}

\begin{example}{Matrice nilpotente}{ex:nilpotente}
  Soit $N=\begin{pmatrix}0&1&0\\0&0&1\\0&0&0\end{pmatrix}$. Calculons~:
  \[
    N^2=\begin{pmatrix}0&0&1\\0&0&0\\0&0&0\end{pmatrix},\qquad
    N^3=\begin{pmatrix}0&0&0\\0&0&0\\0&0&0\end{pmatrix}=0.
  \]
  Donc $N$ est nilpotente d'indice $\nu(N)=3$. On vérifie~:
  $\chi_N(X)=X^3$ et $\mu_N(X)=X^3$.
\end{example}

\begin{proposition}{Suite des noyaux itérés}{prop:noyaux-iteres}
  Soit $f\in\End(E)$ nilpotent d'indice~$\nu$. On a la chaîne
  strictement croissante~:
  \[
    \set{\vzero}=\Ker(f^0)\subsetneq\Ker(f)\subsetneq\Ker(f^2)
    \subsetneq\cdots\subsetneq\Ker(f^\nu)=E.
  \]
  De plus, pour tout $k\geq\nu$, $\Ker(f^k)=E$.
\end{proposition}

\begin{proof}
  La croissance $\Ker(f^k)\subseteq\Ker(f^{k+1})$ est immédiate~: si
  $f^k(\vec{v})=\vzero$, alors
  $f^{k+1}(\vec{v})=f(f^k(\vec{v}))=f(\vzero)=\vzero$.

  Montrons que l'inclusion est stricte pour $k<\nu$. Par définition de
  l'indice de nilpotence, $f^{k}\neq 0$ pour $k<\nu$, donc il existe
  $\vec{w}\in E$ tel que $f^k(\vec{w})\neq\vzero$, c'est-à-dire
  $\vec{w}\notin\Ker(f^k)$. Montrons que si
  $\Ker(f^k)=\Ker(f^{k+1})$, alors $\Ker(f^j)=\Ker(f^k)$ pour tout
  $j\geq k$, par récurrence sur~$j$. Si
  $\vec{v}\in\Ker(f^{k+2})$, alors $f(\vec{v})\in\Ker(f^{k+1})=\Ker(f^k)$,
  donc $f^{k+1}(\vec{v})=f^k(f(\vec{v}))=\vzero$, d'où
  $\vec{v}\in\Ker(f^{k+1})=\Ker(f^k)$. On obtient $\Ker(f^k)=E$ pour
  tout $k<\nu$, ce qui contredit $f^k\neq 0$.

  Pour $k\geq\nu$, $f^\nu=0$ implique $f^k=f^{k-\nu}\circ f^\nu=0$,
  donc $\Ker(f^k)=E$.
\end{proof}


% ============================================================
\section{Sous-espaces invariants}\label{sec:sous-espaces-invariants}
% ============================================================

\begin{definition}{Sous-espace invariant}{def:invariant}
  Soit $f\in\End(E)$. Un sous-espace vectoriel $F$ de~$E$ est dit
  \emph{$f$-invariant}\index{sous-espace invariant} (ou \emph{stable
  par~$f$}) si $f(F)\subseteq F$, c'est-à-dire si pour tout
  $\vec{v}\in F$, on a $f(\vec{v})\in F$.
\end{definition}

\begin{example}{Exemples de sous-espaces invariants}{ex:invariant}
  \begin{enumerate}[label=(\roman*)]
    \item $\set{\vzero}$ et $E$ sont toujours $f$-invariants.
    \item $\Ker(f)$ et $\im(f)$ sont $f$-invariants.
    \item Tout sous-espace propre $E_\lambda(f)=\Ker(f-\lambda\Id)$ est
      $f$-invariant.
    \item Plus généralement, $\Ker(P(f))$ est $f$-invariant pour tout
      polynôme~$P$.\label{item:ker-poly-inv}
  \end{enumerate}
\end{example}

\begin{proof}[Preuve de~\ref{item:ker-poly-inv}]
  Soit $\vec{v}\in\Ker(P(f))$, c'est-à-dire $P(f)(\vec{v})=\vzero$. Comme
  $f$ commute avec $P(f)$ (car $P(f)$ est un polynôme en~$f$), on a~:
  \[
    P(f)(f(\vec{v}))=f(P(f)(\vec{v}))=f(\vzero)=\vzero,
  \]
  donc $f(\vec{v})\in\Ker(P(f))$.
\end{proof}

\begin{proposition}{Restriction à un sous-espace invariant}{prop:restriction-inv}
  Si $F$ est un sous-espace $f$-invariant, alors la restriction
  $f|_F\colon F\to F$ est un endomorphisme de~$F$. De plus~:
  \begin{enumerate}[label=(\roman*)]
    \item Le polynôme minimal de $f|_F$ divise le polynôme minimal
      de~$f$.
    \item Le polynôme caractéristique de $f|_F$ divise le polynôme
      caractéristique de~$f$.
  \end{enumerate}
\end{proposition}

\begin{proof}
  (i) Si $\mu_f(f)=0$, alors pour tout $\vec{v}\in F$,
  $\mu_f(f|_F)(\vec{v})=\mu_f(f)(\vec{v})=\vzero$. Donc $\mu_f$
  annule $f|_F$, et le polynôme minimal de $f|_F$ divise $\mu_f$.

  (ii) Complétons une base $(e_1,\ldots,e_p)$ de~$F$ en une base
  $(e_1,\ldots,e_p,e_{p+1},\ldots,e_n)$ de~$E$. Dans cette base, la
  matrice de~$f$ a la forme
  \[
    \Mat(f)=\begin{pmatrix}A&B\\0&C\end{pmatrix},
  \]
  où $A=\Mat(f|_F)\in\Mat_p(\C)$. Alors
  $\chi_f(X)=\det(XI_n-\Mat(f))=\det(XI_p-A)\cdot\det(XI_{n-p}-C)
  =\chi_{f|_F}(X)\cdot\chi_C(X)$.
\end{proof}


% ============================================================
\section{Sous-espaces caractéristiques}\label{sec:sous-espaces-car}
% ============================================================

Les sous-espaces propres ne suffisent pas en général à décomposer~$E$.
Il faut les \og{}épaissir\fg{} en remplaçant $\Ker(f-\lambda\Id)$ par
$\Ker((f-\lambda\Id)^k)$ pour $k$ assez grand.

\begin{definition}{Sous-espace caractéristique}{def:sous-espace-car}
  Soit $f\in\End(E)$ et $\lambda\in\Spec(f)$. Le \emph{sous-espace
  caractéristique}\index{sous-espace caractéristique} (ou \emph{sous-espace
  propre généralisé}\index{sous-espace propre généralisé}) de $f$
  associé à~$\lambda$ est
  \[
    F_\lambda(f)\deff\Ker\bigl((f-\lambda\Id)^n\bigr)
    =\bigcup_{k=0}^{n}\Ker\bigl((f-\lambda\Id)^k\bigr),
  \]
  où $n=\dim E$.
\end{definition}

\begin{definition}{Vecteur propre généralisé}{def:vp-generalise}
  Un vecteur $\vec{v}\in E\setminus\set{\vzero}$ est un \emph{vecteur
  propre généralisé}\index{vecteur propre généralisé} de $f$ associé à la
  valeur propre~$\lambda$ s'il existe un entier $k\geq 1$ tel que
  $(f-\lambda\Id)^k(\vec{v})=\vzero$.
\end{definition}

\begin{proposition}{Propriétés des sous-espaces caractéristiques}{prop:car-prop}
  Soit $f\in\End(E)$ avec $\Spec(f)=\set{\lambda_1,\ldots,\lambda_r}$
  (valeurs propres distinctes sur~$\C$), et notons $m_i$ la multiplicité
  algébrique de~$\lambda_i$. Alors~:
  \begin{enumerate}[label=(\roman*)]
    \item\label{item:car-inv} Chaque $F_{\lambda_i}(f)$ est un
      sous-espace $f$-invariant.
    \item\label{item:car-contient} $E_{\lambda_i}(f)\subseteq F_{\lambda_i}(f)$.
    \item\label{item:car-dim} $\dim F_{\lambda_i}(f)=m_i$.
    \item\label{item:car-decomp} $E=F_{\lambda_1}(f)\oplus F_{\lambda_2}(f)
      \oplus\cdots\oplus F_{\lambda_r}(f)$.
    \item\label{item:car-nilp} La restriction
      $(f-\lambda_i\Id)|_{F_{\lambda_i}(f)}$ est nilpotente.
  \end{enumerate}
\end{proposition}

\begin{proof}
  \ref{item:car-inv} $F_{\lambda_i}(f)=\Ker((f-\lambda_i\Id)^n)$ est
  le noyau d'un polynôme en~$f$, donc $f$-invariant (cf.\
  
ef{ex:ex:invariant}).

  \ref{item:car-contient} Si $\vec{v}\in E_{\lambda_i}(f)$, alors
  $(f-\lambda_i\Id)(\vec{v})=\vzero$, donc
  $(f-\lambda_i\Id)^n(\vec{v})=\vzero$, c'est-à-dire
  $\vec{v}\in F_{\lambda_i}(f)$.

  \ref{item:car-dim} et \ref{item:car-decomp}~: écrivons le polynôme
  caractéristique sous la forme
  \[
    \chi_f(X)=\prod_{i=1}^{r}(X-\lambda_i)^{m_i},\qquad
    \sum_{i=1}^{r}m_i=n.
  \]
  Posons $P_i(X)=\prod_{j\neq i}(X-\lambda_j)^{m_j}$. Les polynômes
  $P_1,\ldots,P_r$ sont premiers entre eux dans leur ensemble (car les
  $\lambda_i$ sont distincts), donc il existe des polynômes
  $U_1,\ldots,U_r$ tels que
  $\sum_{i=1}^{r}U_i(X)P_i(X)=1$.

  Par le théorème de Cayley-Hamilton, $\chi_f(f)=0$, donc
  $(f-\lambda_i\Id)^{m_i}\circ P_i(f)=0$ pour tout~$i$, ce qui
  entraîne $\im(P_i(f))\subseteq\Ker((f-\lambda_i\Id)^{m_i})\subseteq
  F_{\lambda_i}(f)$.

  De la relation $\sum_i U_i(f)\circ P_i(f)=\Id$, on déduit que pour
  tout $\vec{v}\in E$~:
  \[
    \vec{v}=\sum_{i=1}^{r}U_i(f)(P_i(f)(\vec{v})),
  \]
  et chaque terme $U_i(f)(P_i(f)(\vec{v}))\in F_{\lambda_i}(f)$ (car
  $F_{\lambda_i}(f)$ est $f$-invariant et contient $\im(P_i(f))$).
  Ainsi $E=\sum_i F_{\lambda_i}(f)$.

  Montrons que la somme est directe. Si
  $\vec{v}\in F_{\lambda_i}(f)\cap\sum_{j\neq i}F_{\lambda_j}(f)$,
  alors $(f-\lambda_i\Id)^{m_i}(\vec{v})=\vzero$ car
  $\vec{v}\in F_{\lambda_i}(f)\subseteq\Ker((f-\lambda_i\Id)^{m_i})$.
  Par ailleurs, $\vec{v}=\sum_{j\neq i}\vec{v}_j$ avec
  $\vec{v}_j\in F_{\lambda_j}(f)$, et sur chaque $F_{\lambda_j}(f)$
  ($j\neq i$), l'endomorphisme $f-\lambda_i\Id$ est inversible
  (car $\lambda_j\neq\lambda_i$ implique que $\lambda_i$ n'est pas
  valeur propre de $f|_{F_{\lambda_j}(f)}$). Donc
  $(f-\lambda_i\Id)^{m_i}$ est inversible sur
  $\sum_{j\neq i}F_{\lambda_j}(f)$, ce qui force $\vec{v}=\vzero$.

  La somme directe $E=\bigoplus_i F_{\lambda_i}(f)$ donne
  $n=\sum_i\dim F_{\lambda_i}(f)$. De plus,
  $F_{\lambda_i}(f)\subseteq\Ker((f-\lambda_i\Id)^{m_i})$ (par le
  raisonnement ci-dessus) et l'on sait que
  $\dim\Ker((f-\lambda_i\Id)^{m_i})\leq m_i$ (car le polynôme
  caractéristique de $f|_{F_{\lambda_i}(f)}$ est de degré
  $\dim F_{\lambda_i}(f)$, et $(X-\lambda_i)^{m_i}$ divise $\chi_f$).
  L'égalité $\sum_i\dim F_{\lambda_i}(f)=\sum_i m_i=n$ force
  $\dim F_{\lambda_i}(f)=m_i$.

  \ref{item:car-nilp} Sur $F_{\lambda_i}(f)$, tout vecteur
  $\vec{v}$ vérifie $(f-\lambda_i\Id)^{m_i}(\vec{v})=\vzero$, donc
  $(f-\lambda_i\Id)|_{F_{\lambda_i}(f)}$ est nilpotent (d'indice
  $\leq m_i$).
\end{proof}


% ============================================================
\section{Blocs de Jordan et matrices de Jordan}\label{sec:blocs-jordan}
% ============================================================

\begin{definition}{Bloc de Jordan}{def:bloc-jordan}
  Soit $\lambda\in\C$ et $k\geq 1$ un entier. Le \emph{bloc de Jordan}%
  \index{bloc de Jordan} de taille~$k$ associé à~$\lambda$ est la
  matrice $J_k(\lambda)\in\Mat_k(\C)$ définie par
  \[
    J_k(\lambda)\deff
    \begin{pmatrix}
      \lambda & 1      & 0      & \cdots & 0 \\
      0       &\lambda & 1      & \ddots & \vdots \\
      \vdots  &\ddots  &\ddots  & \ddots & 0 \\
      0       &\cdots  & 0      &\lambda & 1 \\
      0       &\cdots  & \cdots & 0      & \lambda
    \end{pmatrix}
    =\lambda I_k+N_k,
  \]
  où $N_k$ est la matrice nilpotente\index{matrice nilpotente!élémentaire}
  ayant des~$1$ sur la sur-diagonale et des~$0$ partout ailleurs.
\end{definition}

\begin{remark}{Cas particuliers}{rem:bloc-cas-part}
  \begin{itemize}
    \item Pour $k=1$~: $J_1(\lambda)=(\lambda)$, c'est un scalaire.
    \item Un bloc de Jordan $J_k(\lambda)$ avec $\lambda\neq 0$ n'est pas
      nilpotent, mais $J_k(\lambda)-\lambda I_k=N_k$ l'est.
    \item Un bloc $J_k(0)=N_k$ est nilpotent d'indice~$k$.
  \end{itemize}
\end{remark}

\begin{definition}{Matrice de Jordan}{def:matrice-jordan}
  Une \emph{matrice de Jordan}\index{matrice de Jordan} est une matrice
  diagonale par blocs dont chaque bloc diagonal est un bloc de Jordan~:
  \[
    J=\diag\bigl(J_{k_1}(\lambda_1),\,J_{k_2}(\lambda_2),\,\ldots,\,
    J_{k_s}(\lambda_s)\bigr)
    =\begin{pmatrix}
      J_{k_1}(\lambda_1) & & & 0 \\
       & J_{k_2}(\lambda_2) & & \\
       & & \ddots & \\
      0 & & & J_{k_s}(\lambda_s)
    \end{pmatrix}.
  \]
  Les $\lambda_i$ ne sont pas nécessairement distincts, et
  $k_1+k_2+\cdots+k_s=n$.
\end{definition}

\begin{proposition}{Propriétés d'un bloc de Jordan}{prop:bloc-jordan-prop}
  Soit $J_k(\lambda)$ un bloc de Jordan. Alors~:
  \begin{enumerate}[label=(\roman*)]
    \item\label{item:chi-jordan} $\chi_{J_k(\lambda)}(X)=(X-\lambda)^k$.
    \item\label{item:mu-jordan} $\mu_{J_k(\lambda)}(X)=(X-\lambda)^k$.
    \item\label{item:sep-jordan} $\dim E_\lambda(J_k(\lambda))=1$.
    \item\label{item:jordan-diag} $J_k(\lambda)$ est diagonalisable si
      et seulement si $k=1$.
  \end{enumerate}
\end{proposition}

\begin{proof}
  \ref{item:chi-jordan}
  $J_k(\lambda)-XI_k$ est triangulaire supérieure avec $\lambda-X$
  sur la diagonale, donc $\chi_{J_k(\lambda)}(X)=(\lambda-X)^k=(X-\lambda)^k$
  (à signe près, selon la convention; nous utilisons
  $\chi(X)=\det(XI-A)$).

  \ref{item:mu-jordan} On a $(J_k(\lambda)-\lambda I_k)=N_k$ et
  $N_k^k=0$, $N_k^{k-1}\neq 0$. Donc $\mu_{J_k(\lambda)}(X)=(X-\lambda)^k$.

  \ref{item:sep-jordan}
  $\Ker(J_k(\lambda)-\lambda I_k)=\Ker(N_k)$ est engendré par le premier
  vecteur de la base canonique, donc de dimension~$1$.

  \ref{item:jordan-diag} Si $k\geq 2$, la multiplicité algébrique de
  $\lambda$ est~$k$ mais la multiplicité géométrique est~$1$, donc
  $J_k(\lambda)$ n'est pas diagonalisable. Si $k=1$, $J_1(\lambda)=(\lambda)$
  est déjà diagonale.
\end{proof}


% ============================================================
\section{Théorème de la forme normale de Jordan}\label{sec:theoreme-jordan}
% ============================================================

\begin{theorem}{Forme normale de Jordan}{thm:jordan}
  Soit $E$ un $\C$-espace vectoriel de dimension $n\geq 1$ et
  $f\in\End(E)$.\index{forme normale de Jordan!théorème} Alors il existe
  une base $\cB$ de~$E$ dans laquelle la matrice de~$f$ est une matrice
  de Jordan~:
  \[
    \Mat_\cB(f)=\diag\bigl(J_{k_1}(\lambda_1),\,\ldots,\,
    J_{k_s}(\lambda_s)\bigr),
  \]
  où $\lambda_1,\ldots,\lambda_s\in\C$ (non nécessairement distincts) et
  $k_1+\cdots+k_s=n$.

  De plus, cette forme de Jordan est unique à l'ordre des blocs près~:
  les tailles des blocs et les valeurs propres associées sont uniquement
  déterminées par~$f$.
\end{theorem}

De manière équivalente, pour toute matrice $A\in\Mat_n(\C)$, il existe
une matrice inversible $P\in\GL_n(\C)$ telle que
$\inv{P}AP=J$ est une matrice de Jordan.\index{matrice semblable!forme de Jordan}

\begin{proof}[Preuve du 
ef{thm:thm:jordan}]
  La preuve se déroule en plusieurs étapes.

  \medskip\noindent
  \textbf{Étape 1~: Décomposition en sous-espaces caractéristiques.}

  Par la 
ef{prop:prop:car-prop}, on dispose de la décomposition
  \[
    E=F_{\lambda_1}(f)\oplus\cdots\oplus F_{\lambda_r}(f),
  \]
  où $\lambda_1,\ldots,\lambda_r$ sont les valeurs propres distinctes
  de~$f$, et $\dim F_{\lambda_i}(f)=m_i$ (multiplicité algébrique).
  Chaque $F_{\lambda_i}(f)$ est $f$-invariant, et la restriction
  $g_i=(f-\lambda_i\Id)|_{F_{\lambda_i}(f)}$ est nilpotente.

  Il suffit donc de montrer qu'un endomorphisme nilpotent d'un espace de
  dimension finie admet une forme de Jordan (avec $\lambda=0$). En effet,
  si dans une base de $F_{\lambda_i}(f)$, la matrice de $g_i$ est
  $\diag(N_{k_1},\ldots,N_{k_t})$ (blocs nilpotents), alors la matrice
  de $f|_{F_{\lambda_i}(f)}$ dans cette même base est
  $\diag(J_{k_1}(\lambda_i),\ldots,J_{k_t}(\lambda_i))$.

  \medskip\noindent
  \textbf{Étape 2~: Forme de Jordan d'un endomorphisme nilpotent.}

  Soit $g\colon V\to V$ un endomorphisme nilpotent d'un espace vectoriel
  $V$ de dimension $m$, d'indice de nilpotence~$\nu$. On procède par
  récurrence forte sur $m=\dim V$.

  \emph{Cas de base~:} si $m=1$, alors $g=0$ et la matrice de $g$ est
  $(0)=J_1(0)$.

  \emph{Étape de récurrence~:} supposons le résultat vrai pour tout
  espace de dimension $<m$. Puisque $g$ est nilpotent et $m\geq 2$,
  $g$ n'est pas inversible, donc $\Ker(g)\neq\set{\vzero}$ et
  $W=\im(g)\subsetneq V$ est un sous-espace propre de~$V$ (de dimension
  $<m$), stable par~$g$.

  Par hypothèse de récurrence, il existe une base de $W$ dans laquelle
  $g|_W$ est de la forme
  \[
    \diag(N_{p_1},\ldots,N_{p_t})
  \]
  avec $p_1\geq p_2\geq\cdots\geq p_t\geq 1$.

  Pour chaque bloc $N_{p_j}$, notons $(w_1^{(j)},\ldots,w_{p_j}^{(j)})$
  les vecteurs de base correspondants dans~$W$, de sorte que
  \[
    g(w_i^{(j)})=w_{i+1}^{(j)} \quad\text{pour } i<p_j,\qquad
    g(w_{p_j}^{(j)})=\vzero.
  \]
  Le vecteur $w_1^{(j)}\in\im(g)$, donc il existe $v_0^{(j)}\in V$ tel
  que $g(v_0^{(j)})=w_1^{(j)}$. Alors la famille
  $(v_0^{(j)},w_1^{(j)},\ldots,w_{p_j}^{(j)})$ vérifie
  \[
    g(v_0^{(j)})=w_1^{(j)},\quad g(w_1^{(j)})=w_2^{(j)},\quad\ldots,\quad
    g(w_{p_j}^{(j)})=\vzero,
  \]
  ce qui donne un bloc de Jordan $N_{p_j+1}$... mais il faut vérifier
  que les $v_0^{(j)}$ peuvent être choisis de manière à former avec les
  vecteurs de base de $W$ et un complément de $\Ker(g)\cap(\im(g))^\perp$
  une base de~$V$.

  Plus précisément, construisons la base de Jordan directement.
  Choisissons $\vec{v}\in V$ tel que $g^{\nu-1}(\vec{v})\neq\vzero$
  (un tel vecteur existe par définition de~$\nu$). Alors les vecteurs
  \[
    \vec{v},\;g(\vec{v}),\;g^2(\vec{v}),\;\ldots,\;g^{\nu-1}(\vec{v})
  \]
  sont linéairement indépendants (car si
  $\sum_{i=0}^{\nu-1}\alpha_i g^i(\vec{v})=\vzero$, en appliquant
  $g^{\nu-1}$ on obtient $\alpha_0 g^{\nu-1}(\vec{v})=\vzero$, d'où
  $\alpha_0=0$, puis par récurrence $\alpha_i=0$ pour tout~$i$).

  Posons $V_1=\Vect(\vec{v},g(\vec{v}),\ldots,g^{\nu-1}(\vec{v}))$.
  C'est un sous-espace de dimension~$\nu$, $g$-invariant, et dans la base
  $(g^{\nu-1}(\vec{v}),\ldots,g(\vec{v}),\vec{v})$, la matrice de $g|_{V_1}$
  est $N_\nu$.

  Si $V_1=V$, on a terminé. Sinon, il faut trouver un supplémentaire
  $g$-invariant de $V_1$ dans~$V$. Cela se fait en choisissant un
  supplémentaire $V_2$ de $V_1$ dans~$V$ tel que $g(V_2)\subseteq V_2$
  (un tel supplémentaire existe, voir ci-dessous), et on applique
  l'hypothèse de récurrence à $g|_{V_2}$.

  \emph{Existence d'un supplémentaire $g$-invariant.} Nous l'établissons
  par le lemme suivant.
\end{proof}

\begin{lemma}{Supplémentaire invariant pour un nilpotent}{lem:suppl-invariant}
  Soit $g\colon V\to V$ nilpotent et soit $V_1$ un sous-espace
  $g$-invariant de~$V$ ayant une base cyclique
  $(g^{\nu-1}(\vec{v}),\ldots,g(\vec{v}),\vec{v})$ pour un certain
  $\vec{v}\in V$ avec $g^\nu(\vec{v})=\vzero$ et
  $g^{\nu-1}(\vec{v})\neq\vzero$. Alors il existe un sous-espace $V_2$
  de~$V$, $g$-invariant, tel que $V=V_1\oplus V_2$.
\end{lemma}

\begin{proof}
  Nous procédons par récurrence sur $\dim V$. Si $\dim V=\nu$, alors
  $V_1=V$ et $V_2=\set{\vzero}$ convient. Supposons $\dim V>\nu$.

  Comme $g^{\nu-1}(\vec{v})\in\Ker(g)\setminus\set{\vzero}$,
  choisissons un sous-espace $W$ de $\Ker(g)$ tel que
  $\Ker(g)=\Vect(g^{\nu-1}(\vec{v}))\oplus W$.

  \emph{Cas $\nu=1$.} Alors $g(\vec{v})=\vzero$, $V_1=\Vect(\vec{v})$
  et $V_1\subseteq\Ker(g)$. Choisissons un supplémentaire $W'$ de $V_1$
  dans $\Ker(g)$, et un supplémentaire $V_2$ de $\Ker(g)$ dans $V$.
  Alors $V=V_1\oplus(W'\oplus V_2)$. On peut vérifier que
  $g(W'\oplus V_2)\subseteq\Ker(g)$... La construction nécessite plus de
  soin. On peut étendre $W$ en un supplémentaire de $V_1$ et vérifier
  la $g$-invariance par récurrence.

  \emph{Cas $\nu\geq 2$.} Posons $V'=g^{-1}(W+V_1)$. Considérons
  $\bar{V}=V/V_1$ et l'endomorphisme induit $\bar{g}\colon\bar{V}\to\bar{V}$
  (bien défini car $V_1$ est $g$-invariant). Alors $\bar{g}$ est
  nilpotent sur $\bar{V}$ et $\dim\bar{V}<\dim V$. Par un argument de
  relèvement, on construit le supplémentaire voulu.

  La construction complète peut se faire par récurrence sur $\nu$~: on
  pose $\vec{v}'=g(\vec{v})$, $V_1'=\Vect(g^{\nu-2}(\vec{v}'),\ldots,
  g(\vec{v}'),\vec{v}')$, et $\bar{V}=\im(g)$. Alors $g|_{\bar{V}}$
  est nilpotent et $V_1'\subseteq\bar{V}$ est un sous-espace cyclique
  de dimension $\nu-1$. Par hypothèse de récurrence, il existe
  $V_2'\subseteq\bar{V}$ tel que $\bar{V}=V_1'\oplus V_2'$, avec $V_2'$
  $g$-stable.

  Pour chaque vecteur de base $w_j$ de $V_2'$, choisissons un
  antécédent $u_j\in V$ tel que $g(u_j)=w_j$. Les vecteurs $u_j$,
  complétés par les chaînes cycliques qu'ils engendrent et les vecteurs
  de $\Ker(g)\cap V_2'$, forment le supplémentaire $V_2$ cherché (on
  vérifie que $V_2$ est $g$-stable et que $V=V_1\oplus V_2$ par un
  argument de dimension).
\end{proof}

\begin{proof}[Fin de la preuve du 
ef{thm:thm:jordan} --- Unicité]
  \textbf{Étape 3~: Unicité.}

  Soit $\lambda$ une valeur propre de~$f$, et notons
  $k_1\geq k_2\geq\cdots\geq k_t$ les tailles des blocs de Jordan
  associés à~$\lambda$. Posons $g=(f-\lambda\Id)|_{F_\lambda(f)}$, qui
  est nilpotent. Pour tout $j\geq 0$, calculons~:
  \[
    r_j\deff\rg(g^j)=\dim F_\lambda(f)-\dim\Ker(g^j).
  \]
  Or, pour un bloc $N_k$ de taille~$k$, on a
  $\dim\Ker(N_k^j)=\min(j,k)$. Donc
  \[
    \dim\Ker(g^j)=\sum_{i=1}^{t}\min(j,k_i).
  \]
  Les quantités $\dim\Ker(g^j)$ sont des invariants de~$f$
  (elles ne dépendent pas de la base choisie). À partir de la suite
  $(\dim\Ker(g^j))_{j\geq 0}$, on retrouve les tailles des blocs~:
  le nombre de blocs de taille $\geq j$ est
  $\dim\Ker(g^j)-\dim\Ker(g^{j-1})$, et le nombre de blocs de
  taille exactement~$j$ est
  \[
    \bigl(\dim\Ker(g^j)-\dim\Ker(g^{j-1})\bigr)
    -\bigl(\dim\Ker(g^{j+1})-\dim\Ker(g^j)\bigr)
    =2\dim\Ker(g^j)-\dim\Ker(g^{j-1})-\dim\Ker(g^{j+1}).
  \]
  Les tailles de blocs sont donc uniquement déterminées.
\end{proof}


% ============================================================
\section{Calcul de la forme de Jordan}\label{sec:calcul-jordan}
% ============================================================

Pour calculer la forme de Jordan d'une matrice $A\in\Mat_n(\C)$, on suit
l'algorithme suivant.\index{forme de Jordan!calcul}

\begin{enumerate}
  \item \textbf{Trouver les valeurs propres~:} calculer le polynôme
    caractéristique $\chi_A(X)=\det(XI_n-A)$ et le factoriser sur~$\C$~:
    \[
      \chi_A(X)=\prod_{i=1}^{r}(X-\lambda_i)^{m_i}.
    \]
    La multiplicité algébrique de $\lambda_i$ est~$m_i$.

  \item \textbf{Pour chaque valeur propre $\lambda_i$, déterminer les
    tailles des blocs~:} calculer successivement~:
    \[
      d_j^{(i)}=\dim\Ker\bigl((A-\lambda_i I_n)^j\bigr),
      \qquad j=1,2,3,\ldots
    \]
    jusqu'à atteindre $d_j^{(i)}=m_i$. On pose $d_0^{(i)}=0$.

  \item \textbf{Déterminer le nombre de blocs de chaque taille~:}
    le nombre de blocs de Jordan de taille~$j$ associés à $\lambda_i$
    est
    \[
      b_j^{(i)}=d_j^{(i)}-2d_{j-1}^{(i)}+d_{j-2}^{(i)},
    \]
    où $d_{-1}^{(i)}=0$ par convention. De manière équivalente,
    le nombre total de blocs associés à $\lambda_i$ est
    $d_1^{(i)}$ (la multiplicité géométrique), et le nombre de blocs de
    taille $\geq j$ est $d_j^{(i)}-d_{j-1}^{(i)}$.

  \item \textbf{Construire la matrice de Jordan~:} assembler les
    blocs $J_{k}(\lambda_i)$ selon les tailles déterminées.
\end{enumerate}

\begin{remark}{Vérifications}{rem:verif-jordan}
  On vérifie toujours que~:
  \begin{itemize}
    \item la somme des tailles de tous les blocs vaut~$n$,
    \item pour chaque $\lambda_i$, la somme des tailles des blocs
      associés vaut~$m_i$,
    \item le nombre de blocs de Jordan associés à $\lambda_i$ égale
      la multiplicité géométrique $\dim E_{\lambda_i}(A)$.
  \end{itemize}
\end{remark}


% ============================================================
\section{Exemples détaillés}\label{sec:exemples-jordan}
% ============================================================

\begin{example}{Matrice $3\times 3$}{ex:jordan-3x3}
  Soit $A=\begin{pmatrix}2&1&0\\0&2&1\\0&0&2\end{pmatrix}$.

  \textbf{Étape 1~:} $A$ est triangulaire, donc ses valeurs propres sont
  les éléments diagonaux~: $\lambda=2$ de multiplicité algébrique $m=3$.

  \textbf{Étape 2~:} Posons $B=A-2I_3=\begin{pmatrix}0&1&0\\0&0&1\\0&0&0\end{pmatrix}$.
  \begin{align*}
    d_1&=\dim\Ker(B)=1, \\
    d_2&=\dim\Ker(B^2)=2,\quad\text{car }
      B^2=\begin{pmatrix}0&0&1\\0&0&0\\0&0&0\end{pmatrix}, \\
    d_3&=\dim\Ker(B^3)=3,\quad\text{car }B^3=0.
  \end{align*}

  \textbf{Étape 3~:} Nombre de blocs de taille~$j$~:
  \begin{itemize}
    \item $j=1$~: $b_1=d_1-2d_0+d_{-1}=1-0+0=1$... Utilisons
      plutôt la formule des différences successives.
  \end{itemize}
  Nombre de blocs de taille $\geq j$~:
  \begin{align*}
    \Delta_1&=d_1-d_0=1-0=1, \\
    \Delta_2&=d_2-d_1=2-1=1, \\
    \Delta_3&=d_3-d_2=3-2=1.
  \end{align*}
  Nombre de blocs de taille exactement~$j$~: $\Delta_j-\Delta_{j+1}$~:
  \begin{itemize}
    \item taille $1$~: $\Delta_1-\Delta_2=1-1=0$,
    \item taille $2$~: $\Delta_2-\Delta_3=1-1=0$,
    \item taille $3$~: $\Delta_3-0=1$.
  \end{itemize}
  Il y a un seul bloc de taille~$3$.

  \textbf{Conclusion~:} La forme de Jordan de~$A$ est
  \[
    J=J_3(2)=\begin{pmatrix}2&1&0\\0&2&1\\0&0&2\end{pmatrix}=A.
  \]
  La matrice $A$ est déjà sous forme de Jordan.
\end{example}

\begin{example}{Matrice $4\times 4$}{ex:jordan-4x4}
  Soit $A=\begin{pmatrix}3&1&0&0\\0&3&0&0\\0&0&3&1\\0&0&0&3\end{pmatrix}$.

  \textbf{Étape 1~:} $\chi_A(X)=(X-3)^4$, donc $\lambda=3$ de multiplicité~$4$.

  \textbf{Étape 2~:} $B=A-3I_4=\begin{pmatrix}0&1&0&0\\0&0&0&0\\0&0&0&1\\0&0&0&0\end{pmatrix}$.
  \begin{align*}
    d_1&=\dim\Ker(B)=2\quad(\text{les colonnes 1 et 3 sont nulles}), \\
    d_2&=\dim\Ker(B^2)=4\quad(\text{car }B^2=0).
  \end{align*}

  \textbf{Étape 3~:}
  $\Delta_1=d_1-d_0=2$, $\Delta_2=d_2-d_1=2$.
  Blocs de taille exactement $1$~: $\Delta_1-\Delta_2=0$.
  Blocs de taille exactement $2$~: $\Delta_2=2$.

  \textbf{Conclusion~:} La forme de Jordan est
  \[
    J=\diag(J_2(3),J_2(3))
    =\begin{pmatrix}3&1&0&0\\0&3&0&0\\0&0&3&1\\0&0&0&3\end{pmatrix}=A.
  \]
\end{example}

\begin{example}{Matrice $4\times 4$ avec deux valeurs propres}{ex:jordan-4x4-deux-vp}
  Soit $A=\begin{pmatrix}1&1&0&0\\0&1&0&0\\0&0&2&1\\0&0&0&2\end{pmatrix}$.

  Le polynôme caractéristique est $\chi_A(X)=(X-1)^2(X-2)^2$.

  \emph{Pour $\lambda_1=1$~:} $B_1=A-I_4$. On trouve
  $d_1=\dim\Ker(B_1)=1$, $d_2=\dim\Ker(B_1^2)=2$. Donc un seul
  bloc $J_2(1)$.

  \emph{Pour $\lambda_2=2$~:} $B_2=A-2I_4$. On trouve
  $d_1=\dim\Ker(B_2)=1$, $d_2=\dim\Ker(B_2^2)=2$. Donc un seul
  bloc $J_2(2)$.

  La forme de Jordan est~:
  \[
    J=\diag(J_2(1),J_2(2))
    =\begin{pmatrix}1&1&0&0\\0&1&0&0\\0&0&2&1\\0&0&0&2\end{pmatrix}.
  \]
\end{example}

\begin{example}{Matrice non triviale}{ex:jordan-nontrivial}
  Soit $A=\begin{pmatrix}5&4&2&1\\0&1&-1&-1\\-1&-1&3&0\\1&1&-1&2\end{pmatrix}$.

  On calcule $\chi_A(X)=(X-1)(X-2)(X-4)^2$. Les valeurs propres sont
  $\lambda_1=1$, $\lambda_2=2$, $\lambda_3=4$.

  Pour $\lambda_1=1$ et $\lambda_2=2$, la multiplicité algébrique
  est~$1$, donc le bloc associé est $J_1(1)=(1)$ et $J_1(2)=(2)$.

  Pour $\lambda_3=4$ de multiplicité~$2$~: on calcule
  $d_1=\dim\Ker(A-4I_4)$. Si $d_1=2$, il y a deux blocs $J_1(4)$
  (diagonalisable pour cette valeur propre). Si $d_1=1$, il y a un
  bloc $J_2(4)$.

  On calcule
  $A-4I_4=\begin{pmatrix}1&4&2&1\\0&-3&-1&-1\\-1&-1&-1&0\\1&1&-1&-2\end{pmatrix}$.
  Par réduction, on trouve $\rg(A-4I_4)=3$, donc $d_1=1$ et la forme
  de Jordan est
  \[
    J=\diag\bigl(J_1(1),\,J_1(2),\,J_2(4)\bigr)
    =\begin{pmatrix}1&0&0&0\\0&2&0&0\\0&0&4&1\\0&0&0&4\end{pmatrix}.
  \]
\end{example}


% ============================================================
\section{Forme de Jordan et polynôme minimal}\label{sec:jordan-polymin}
% ============================================================

\begin{theorem}{Polynôme minimal et blocs de Jordan}{thm:polymin-jordan}
  Soit $f\in\End(E)$ de forme de Jordan
  $J=\diag(J_{k_1}(\lambda_1),\ldots,J_{k_s}(\lambda_s))$, et soient
  $\mu_1,\ldots,\mu_r$ les valeurs propres distinctes de~$f$.
  Pour chaque $\mu_i$, notons $\nu_i$ la taille du plus grand bloc de
  Jordan associé à~$\mu_i$.\index{polynôme minimal!et forme de Jordan}
  Alors le polynôme minimal de~$f$ est
  \[
    \mu_f(X)=\prod_{i=1}^{r}(X-\mu_i)^{\nu_i}.
  \]
\end{theorem}

\begin{proof}
  Posons $Q(X)=\prod_{i=1}^{r}(X-\mu_i)^{\nu_i}$. Montrons que
  $Q(f)=0$ et que $\mu_f$ divise~$Q$ et $Q$ divise~$\mu_f$.

  Pour chaque bloc $J_k(\mu_i)$, on a
  $(J_k(\mu_i)-\mu_i I_k)^{\nu_i}=N_k^{\nu_i}=0$ car $k\leq\nu_i$.
  De plus, les facteurs $(X-\mu_j)^{\nu_j}$ pour $j\neq i$ donnent
  des matrices inversibles sur le bloc $J_k(\mu_i)$. Donc
  $Q(J_k(\mu_i))=0$ pour chaque bloc, et $Q(J)=0$, d'où $Q(f)=0$.
  Ainsi $\mu_f\mid Q$.

  Inversement, $\mu_f(f)=0$ implique $\mu_f(J_k(\mu_i))=0$ pour chaque
  bloc. En écrivant $\mu_f(X)=\prod_i(X-\mu_i)^{a_i}$, la condition
  $\mu_f(J_k(\mu_i))=0$ impose $N_k^{a_i}\cdot C=0$ où $C$ est
  inversible (produit des $(J_k(\mu_i)-\mu_j I_k)^{a_j}$ pour $j\neq i$).
  Donc $N_k^{a_i}=0$, ce qui requiert $a_i\geq k$. Prenant le plus
  grand bloc ($k=\nu_i$), on obtient $a_i\geq\nu_i$, d'où $Q\mid\mu_f$.

  Finalement $\mu_f=Q$.
\end{proof}

\begin{corollary}{Diagonalisabilité et blocs}{cor:diag-blocs}
  $f$ est diagonalisable si et seulement si tous les blocs de Jordan
  sont de taille~$1$, c'est-à-dire si et seulement si le polynôme
  minimal est à racines simples.\index{diagonalisabilité!et blocs de Jordan}
\end{corollary}

\begin{proof}
  $f$ est diagonalisable si et seulement si $\mu_f$ est scindé à
  racines simples, c'est-à-dire $\nu_i=1$ pour tout~$i$, ce qui
  signifie que tous les blocs de Jordan sont de taille~$1$.
\end{proof}


% ============================================================
\section{Cayley-Hamilton via la forme de Jordan}\label{sec:cayley-hamilton-jordan}
% ============================================================

\begin{theorem}{Théorème de Cayley-Hamilton (nouvelle preuve)}{thm:cayley-hamilton-jordan}
  \index{théorème de Cayley-Hamilton!preuve par Jordan}
  Soit $A\in\Mat_n(\C)$. Alors $\chi_A(A)=0$.
\end{theorem}

\begin{proof}
  Par le 
ef{thm:thm:jordan}, il existe $P\in\GL_n(\C)$ telle que
  $\inv{P}AP=J=\diag(J_{k_1}(\lambda_1),\ldots,J_{k_s}(\lambda_s))$.

  Comme $\chi_A=\chi_J$, on a~:
  \[
    \chi_A(A)=P\chi_J(J)\inv{P}.
  \]
  Il suffit de montrer $\chi_J(J)=0$. Puisque
  $\chi_J(X)=\prod_{j=1}^{s}(X-\lambda_j)^{k_j}$ et que $J$ est
  diagonale par blocs~:
  \[
    \chi_J(J)=\diag\bigl(\chi_J(J_{k_1}(\lambda_1)),\ldots,
    \chi_J(J_{k_s}(\lambda_s))\bigr).
  \]
  Pour chaque bloc $J_{k_j}(\lambda_j)$~:
  \[
    \chi_J(J_{k_j}(\lambda_j))
    =\prod_{i=1}^{s}(J_{k_j}(\lambda_j)-\lambda_i I_{k_j})^{k_i}.
  \]
  Le facteur correspondant à $i=j$ donne
  $(J_{k_j}(\lambda_j)-\lambda_j I_{k_j})^{k_j}=N_{k_j}^{k_j}=0$.
  Donc l'ensemble du produit est nul, et $\chi_J(J_{k_j}(\lambda_j))=0$.

  Ceci étant vrai pour chaque bloc, $\chi_J(J)=0$ et $\chi_A(A)=0$.
\end{proof}

\begin{remark}{Avantage de cette preuve}{rem:cayley-jordan}
  Cette preuve est conceptuellement très claire~: elle réduit le
  théorème de Cayley-Hamilton au cas des blocs de Jordan, pour lesquels
  la vérification est immédiate. Elle requiert toutefois le théorème de
  Jordan, qui est plus profond que Cayley-Hamilton.
\end{remark}


% ============================================================
\section{Exponentielle de matrice}\label{sec:exp-matrice}
% ============================================================

\begin{definition}{Exponentielle de matrice}{def:exp-matrice}
  Pour toute matrice $A\in\Mat_n(\C)$, l'\emph{exponentielle de matrice}%
  \index{exponentielle de matrice} est définie par la série~:
  \[
    e^A\deff\exp(A)\deff\sum_{k=0}^{\infty}\frac{A^k}{k!}
    =I_n+A+\frac{A^2}{2!}+\frac{A^3}{3!}+\cdots
  \]
  Cette série converge absolument pour toute matrice~$A$.
\end{definition}

\begin{proposition}{Propriétés de l'exponentielle}{prop:exp-prop}
  \begin{enumerate}[label=(\roman*)]
    \item $e^{0}=I_n$.
    \item Si $AB=BA$, alors $e^{A+B}=e^Ae^B$.
    \item $e^A$ est toujours inversible, et $(e^A)^{-1}=e^{-A}$.
    \item Si $A=\inv{P}BP$, alors $e^A=\inv{P}e^BP$.
    \item $\det(e^A)=e^{\tr(A)}$.
  \end{enumerate}
\end{proposition}

\begin{proof}
  (i) Immédiat~: $e^0=\sum_{k\geq 0}\frac{0^k}{k!}=I_n$.

  (ii) Lorsque $AB=BA$, le binôme de Newton s'applique~:
  $(A+B)^k=\sum_{j=0}^{k}\binom{k}{j}A^jB^{k-j}$. On vérifie alors
  que le produit de Cauchy des séries de $e^A$ et $e^B$ donne $e^{A+B}$.

  (iii) Comme $A$ et $-A$ commutent, $e^Ae^{-A}=e^{A-A}=e^0=I_n$.

  (iv) $e^A=\sum_k\frac{A^k}{k!}=\sum_k\frac{(\inv{P}BP)^k}{k!}
  =\sum_k\frac{\inv{P}B^kP}{k!}=\inv{P}\bigl(\sum_k\frac{B^k}{k!}\bigr)P
  =\inv{P}e^BP$.

  (v) Par la propriété (iv), il suffit de montrer le résultat pour les
  matrices triangulaires (ou de Jordan). Pour $J=\diag(J_{k_1}(\lambda_1),
  \ldots)$, $\det(e^J)=\prod_j e^{\lambda_j k_j}$ (à préciser via les
  blocs), et $\tr(J)=\sum_j\lambda_j k_j$, d'où le résultat. Plus
  rigoureusement, $e^{J_k(\lambda)}$ est triangulaire supérieure avec
  $e^\lambda$ sur la diagonale, donc $\det(e^{J_k(\lambda)})=e^{k\lambda}$
  et $\tr(J_k(\lambda))=k\lambda$.
\end{proof}

\begin{proposition}{Exponentielle d'un bloc de Jordan}{prop:exp-bloc}
  L'exponentielle du bloc $tJ_k(\lambda)$ (où $t\in\C$) est~:
  \[
    e^{tJ_k(\lambda)}=e^{\lambda t}
    \begin{pmatrix}
      1      & t      & \frac{t^2}{2!} & \cdots & \frac{t^{k-1}}{(k-1)!} \\
      0      & 1      & t              & \cdots & \frac{t^{k-2}}{(k-2)!} \\
      \vdots & \ddots & \ddots         & \ddots & \vdots \\
      0      & \cdots & 0              & 1      & t \\
      0      & \cdots & \cdots         & 0      & 1
    \end{pmatrix}.
  \]
\end{proposition}

\begin{proof}
  On écrit $tJ_k(\lambda)=\lambda t I_k+tN_k$, où $\lambda t I_k$ et
  $tN_k$ commutent. Donc~:
  \[
    e^{tJ_k(\lambda)}=e^{\lambda t I_k}\cdot e^{tN_k}
    =e^{\lambda t}I_k\cdot e^{tN_k}.
  \]
  Or $N_k$ est nilpotent d'indice~$k$, donc la série exponentielle est
  finie~:
  \[
    e^{tN_k}=\sum_{j=0}^{k-1}\frac{t^j}{j!}N_k^j.
  \]
  La matrice $N_k^j$ a des~$1$ sur la $j$-ième sur-diagonale, ce qui
  donne le résultat annoncé.
\end{proof}

\begin{example}{Exponentielle d'une matrice $2\times 2$}{ex:exp-2x2}
  Soit $A=\begin{pmatrix}3&1\\0&3\end{pmatrix}=J_2(3)$. Alors~:
  \[
    e^{tA}=e^{3t}\begin{pmatrix}1&t\\0&1\end{pmatrix}
    =\begin{pmatrix}e^{3t}&te^{3t}\\0&e^{3t}\end{pmatrix}.
  \]
\end{example}

\begin{example}{Exponentielle d'une matrice $3\times 3$}{ex:exp-3x3}
  Soit $A=\begin{pmatrix}2&1&0\\0&2&1\\0&0&2\end{pmatrix}=J_3(2)$. Alors~:
  \[
    e^{tA}=e^{2t}\begin{pmatrix}1&t&\frac{t^2}{2}\\0&1&t\\0&0&1\end{pmatrix}
    =\begin{pmatrix}e^{2t}&te^{2t}&\frac{t^2}{2}e^{2t}\\
    0&e^{2t}&te^{2t}\\0&0&e^{2t}\end{pmatrix}.
  \]
\end{example}


% ============================================================
\section{Applications}\label{sec:applications-jordan}
% ============================================================

\subsection{Systèmes d'équations différentielles linéaires}%
\label{subsec:edo-jordan}

Considérons le système différentiel linéaire à coefficients constants
\index{système différentiel!non diagonalisable}
\begin{equation}\label{eq:systeme-diff}
  X'(t)=AX(t),\qquad X(0)=X_0,
\end{equation}
où $A\in\Mat_n(\C)$ et $X(t)\in\C^n$.

\begin{theorem}{Solution du système différentiel}{thm:sol-edo}
  L'unique solution du problème de Cauchy~\eqref{eq:systeme-diff} est
  \[
    X(t)=e^{tA}X_0.
  \]
\end{theorem}

\begin{proof}
  \emph{Existence~:} posons $X(t)=e^{tA}X_0$. Alors
  $X'(t)=Ae^{tA}X_0=AX(t)$ (car $\frac{d}{dt}e^{tA}=Ae^{tA}$, ce qui
  se vérifie terme à terme sur la série) et $X(0)=e^0X_0=X_0$.

  \emph{Unicité~:} si $Y(t)$ est une autre solution, posons
  $Z(t)=e^{-tA}Y(t)$. Alors
  $Z'(t)=-Ae^{-tA}Y(t)+e^{-tA}Y'(t)=-Ae^{-tA}Y(t)+e^{-tA}AY(t)=0$.
  Donc $Z$ est constante~: $Z(t)=Z(0)=X_0$, d'où $Y(t)=e^{tA}X_0$.
\end{proof}

\begin{example}{Système $2\times 2$ non diagonalisable}{ex:edo-2x2}
  Résoudre $X'=AX$ avec $A=\begin{pmatrix}3&1\\0&3\end{pmatrix}$ et
  $X(0)=\begin{pmatrix}x_0\\y_0\end{pmatrix}$.

  On a $e^{tA}=\begin{pmatrix}e^{3t}&te^{3t}\\0&e^{3t}\end{pmatrix}$,
  donc~:
  \[
    X(t)=\begin{pmatrix}e^{3t}&te^{3t}\\0&e^{3t}\end{pmatrix}
    \begin{pmatrix}x_0\\y_0\end{pmatrix}
    =\begin{pmatrix}(x_0+y_0t)e^{3t}\\y_0e^{3t}\end{pmatrix}.
  \]
  On retrouve le phénomène caractéristique~: la présence du terme $te^{3t}$
  (croissance polynomiale multipliée par l'exponentielle) est la signature
  d'un bloc de Jordan non trivial.
\end{example}

\begin{example}{Système $3\times 3$}{ex:edo-3x3}
  Résoudre $X'=AX$ avec
  $A=\begin{pmatrix}1&0&0\\0&2&1\\0&0&2\end{pmatrix}$.

  La matrice $A$ est déjà sous forme de Jordan~:
  $A=\diag(J_1(1),J_2(2))$. Donc~:
  \[
    e^{tA}=\diag\bigl(e^t,\,e^{2t}\begin{pmatrix}1&t\\0&1\end{pmatrix}\bigr)
    =\begin{pmatrix}e^t&0&0\\0&e^{2t}&te^{2t}\\0&0&e^{2t}\end{pmatrix}.
  \]
  La solution générale est~:
  \[
    X(t)=\begin{pmatrix}c_1e^t\\c_2e^{2t}+c_3te^{2t}\\c_3e^{2t}\end{pmatrix},
  \]
  où $c_1,c_2,c_3$ sont déterminés par la condition initiale.
\end{example}


\subsection{Fonctions de matrices}\label{subsec:fonctions-matrices}

La forme de Jordan permet de définir et calculer des fonctions de
matrices.\index{fonctions de matrices}

\begin{definition}{Fonction d'une matrice}{def:fonction-matrice}
  Soit $f$ une fonction analytique au voisinage de chaque valeur propre
  de $A$. Si $A=P J\inv{P}$ avec $J=\diag(J_{k_1}(\lambda_1),\ldots,
  J_{k_s}(\lambda_s))$, on définit~:
  \[
    f(A)\deff P\,\diag\bigl(f(J_{k_1}(\lambda_1)),\ldots,
    f(J_{k_s}(\lambda_s))\bigr)\,\inv{P},
  \]
  où, pour un bloc de Jordan~:
  \[
    f(J_k(\lambda))=\begin{pmatrix}
      f(\lambda)   & f'(\lambda)  & \frac{f''(\lambda)}{2!}
        & \cdots & \frac{f^{(k-1)}(\lambda)}{(k-1)!} \\
      0 & f(\lambda) & f'(\lambda) & \cdots
        & \frac{f^{(k-2)}(\lambda)}{(k-2)!} \\
      \vdots & \ddots & \ddots & \ddots & \vdots \\
      0 & \cdots & 0 & f(\lambda) & f'(\lambda) \\
      0 & \cdots & \cdots & 0 & f(\lambda)
    \end{pmatrix}.
  \]
\end{definition}

\begin{example}{Racine carrée d'une matrice}{ex:racine-matrice}
  Soit $A=\begin{pmatrix}4&1\\0&4\end{pmatrix}=J_2(4)$. Avec
  $f(x)=\sqrt{x}$, $f'(x)=\frac{1}{2\sqrt{x}}$~:
  \[
    \sqrt{A}=\begin{pmatrix}\sqrt{4}&\frac{1}{2\sqrt{4}}\\0&\sqrt{4}\end{pmatrix}
    =\begin{pmatrix}2&\frac{1}{4}\\0&2\end{pmatrix}.
  \]
  Vérifions~: $\bigl(\sqrt{A}\bigr)^2=
  \begin{pmatrix}4&1\\0&4\end{pmatrix}=A$.
\end{example}


\subsection{Puissances d'une matrice}\label{subsec:puissances}

\begin{proposition}{Puissances d'un bloc de Jordan}{prop:puissance-jordan}
  Pour $m\geq 0$~:\index{puissances de matrices!bloc de Jordan}
  \[
    \bigl(J_k(\lambda)\bigr)^m=\begin{pmatrix}
      \lambda^m & \binom{m}{1}\lambda^{m-1} & \binom{m}{2}\lambda^{m-2}
        & \cdots & \binom{m}{k-1}\lambda^{m-k+1} \\
      0 & \lambda^m & \binom{m}{1}\lambda^{m-1} & \cdots
        & \binom{m}{k-2}\lambda^{m-k+2} \\
      \vdots & \ddots & \ddots & \ddots & \vdots \\
      0 & \cdots & 0 & \lambda^m & \binom{m}{1}\lambda^{m-1} \\
      0 & \cdots & \cdots & 0 & \lambda^m
    \end{pmatrix},
  \]
  où $\binom{m}{j}=\frac{m!}{j!(m-j)!}$ pour $j\leq m$ et
  $\binom{m}{j}=0$ pour $j>m$.
\end{proposition}

\begin{proof}
  On écrit $J_k(\lambda)=\lambda I_k+N_k$. Puisque $\lambda I_k$ et
  $N_k$ commutent, le binôme de Newton donne~:
  \[
    J_k(\lambda)^m=\sum_{j=0}^{k-1}\binom{m}{j}\lambda^{m-j}N_k^j.
  \]
  Le coefficient en position $(i,i+j)$ de $N_k^j$ est~$1$ (pour
  $1\leq i\leq k-j$), ce qui donne la formule annoncée.
\end{proof}


% ============================================================
\section{Exercices}\label{sec:exos-jordan}
% ============================================================

\begin{exercise}{Endomorphismes nilpotents~: propriétés de base}{exo:nilp-base}
  \basic
  Soit $f\in\End(E)$ nilpotent.
  \begin{enumerate}[label=(\alph*)]
    \item Montrer que $\Id+f$ est inversible. \emph{Indication~: utiliser
      la série géométrique finie $(\Id+f)^{-1}=\sum_{k=0}^{\nu-1}(-1)^kf^k$.}
    \item Montrer que si $g$ commute avec $f$ et est nilpotent, alors
      $f+g$ est nilpotent.
    \item Montrer que $\tr(f^k)=0$ pour tout $k\geq 1$.
  \end{enumerate}
\end{exercise}

\begin{exercise}{Classification des nilpotents en dimension 3}{exo:nilp-dim3}
  \basic
  Déterminer toutes les formes de Jordan possibles pour un endomorphisme
  nilpotent d'un espace de dimension~$3$. Pour chaque cas, préciser
  l'indice de nilpotence, le polynôme minimal et la dimension du noyau.
\end{exercise}

\begin{exercise}{Calcul de forme de Jordan}{exo:jordan-calcul-1}
  \intermediate
  Déterminer la forme de Jordan de chacune des matrices suivantes~:
  \begin{enumerate}[label=(\alph*)]
    \item $A=\begin{pmatrix}2&1&0\\0&2&0\\0&0&3\end{pmatrix}$.
    \item $B=\begin{pmatrix}0&1&0\\0&0&1\\0&0&0\end{pmatrix}$.
    \item $C=\begin{pmatrix}1&1&0&0\\0&1&1&0\\0&0&1&0\\0&0&0&2\end{pmatrix}$.
  \end{enumerate}
\end{exercise}

\begin{exercise}{Forme de Jordan d'une matrice $4\times 4$}{exo:jordan-4x4}
  \intermediate
  Soit $A=\begin{pmatrix}2&0&0&0\\1&2&0&0\\0&1&2&0\\0&0&0&2\end{pmatrix}$.
  \begin{enumerate}[label=(\alph*)]
    \item Calculer le polynôme caractéristique de~$A$.
    \item Calculer $\dim\Ker(A-2I_4)$, $\dim\Ker(A-2I_4)^2$,
      $\dim\Ker(A-2I_4)^3$.
    \item En déduire la forme de Jordan de~$A$.
    \item Déterminer le polynôme minimal de~$A$.
  \end{enumerate}
\end{exercise}

\begin{exercise}{Exponentielle de matrice}{exo:exp-mat}
  \intermediate
  Calculer $e^{tA}$ pour~:
  \begin{enumerate}[label=(\alph*)]
    \item $A=\begin{pmatrix}0&1\\0&0\end{pmatrix}$.
    \item $A=\begin{pmatrix}1&1\\0&1\end{pmatrix}$.
    \item $A=\begin{pmatrix}0&1&0\\0&0&1\\0&0&0\end{pmatrix}$.
    \item $A=\begin{pmatrix}2&1&0\\0&2&0\\0&0&3\end{pmatrix}$.
  \end{enumerate}
\end{exercise}

\begin{exercise}{Système différentiel}{exo:ode-jordan}
  \intermediate
  Résoudre le système différentiel $X'=AX$, $X(0)=X_0$, où
  \[
    A=\begin{pmatrix}1&1&0\\0&1&1\\0&0&1\end{pmatrix},\qquad
    X_0=\begin{pmatrix}1\\0\\1\end{pmatrix}.
  \]
\end{exercise}

\begin{exercise}{Polynôme minimal et forme de Jordan}{exo:polymin-jordan}
  \intermediate
  Pour chacun des cas suivants, donner toutes les formes de Jordan
  possibles (à l'ordre des blocs près) pour une matrice $A\in\Mat_5(\C)$~:
  \begin{enumerate}[label=(\alph*)]
    \item $\chi_A(X)=(X-1)^3(X-2)^2$ et $\mu_A(X)=(X-1)^2(X-2)$.
    \item $\chi_A(X)=(X-3)^5$ et $\mu_A(X)=(X-3)^3$.
    \item $\chi_A(X)=(X-1)^5$ et $\mu_A(X)=(X-1)^2$.
  \end{enumerate}
\end{exercise}

\begin{exercise}{Matrice semblable à son bloc de Jordan}{exo:matrice-passage}
  \challenging
  Soit $A=\begin{pmatrix}5&-4\\1&1\end{pmatrix}$.
  \begin{enumerate}[label=(\alph*)]
    \item Montrer que $A$ n'est pas diagonalisable sur~$\C$.
    \item Déterminer sa forme de Jordan~$J$.
    \item Trouver une matrice $P$ telle que $A=PJ\inv{P}$.
    \item Calculer $A^m$ pour tout $m\geq 0$.
    \item Calculer $e^{tA}$.
  \end{enumerate}
\end{exercise}

\begin{exercise}{Commutant d'un bloc de Jordan}{exo:commutant}
  \challenging
  Soit $J=J_n(\lambda)$ un bloc de Jordan de taille~$n$.
  \begin{enumerate}[label=(\alph*)]
    \item Montrer qu'une matrice $B\in\Mat_n(\C)$ commute avec $J$ si
      et seulement si $B$ est un polynôme (de degré $\leq n-1$) en la
      matrice nilpotente $N_n=J-\lambda I_n$.
    \item En déduire que le commutant de $J$ est de dimension~$n$.
    \item Montrer que deux blocs de Jordan $J_n(\lambda)$ et $J_m(\mu)$
      (avec $(n,\lambda)\neq(m,\mu)$) ne sont pas semblables.
  \end{enumerate}
\end{exercise}

\begin{exercise}{Logarithme de matrice}{exo:log-matrice}
  \challenging
  Soit $A\in\Mat_n(\C)$ dont toutes les valeurs propres ont une partie
  réelle strictement positive.
  \begin{enumerate}[label=(\alph*)]
    \item Montrer qu'il existe une matrice $B\in\Mat_n(\C)$ telle que
      $e^B=A$. \emph{Indication~: utiliser la forme de Jordan et le
      logarithme complexe.}
    \item Calculer explicitement un tel $B$ pour
      $A=\begin{pmatrix}e&1\\0&e\end{pmatrix}$.
  \end{enumerate}
\end{exercise}

\begin{exercise}{Décomposition de Dunford}{exo:dunford}
  \intermediate\index{décomposition de Dunford}
  Soit $f\in\End(E)$ avec $E$ un $\C$-espace vectoriel de dimension
  finie.
  \begin{enumerate}[label=(\alph*)]
    \item Montrer qu'il existe un unique couple $(d,n)\in\End(E)^2$ tel
      que~:
      \begin{itemize}
        \item $f=d+n$,
        \item $d$ est diagonalisable,
        \item $n$ est nilpotent,
        \item $d\circ n=n\circ d$.
      \end{itemize}
      \emph{Indication~: utiliser la forme de Jordan pour construire $d$
      et~$n$.}
    \item Montrer que $d$ et $n$ sont des polynômes en~$f$.
    \item Calculer la décomposition de Dunford de
      $A=\begin{pmatrix}3&1&0\\0&3&1\\0&0&3\end{pmatrix}$.
  \end{enumerate}
\end{exercise}

\begin{exercise}{Résolution d'une récurrence linéaire}{exo:recurrence}
  \intermediate
  Considérons la récurrence linéaire
  \[
    u_{n+2}-4u_{n+1}+4u_n=0,\qquad u_0=1,\quad u_1=3.
  \]
  \begin{enumerate}[label=(\alph*)]
    \item Écrire cette récurrence sous forme matricielle
      $X_{n+1}=AX_n$ avec $X_n=\begin{pmatrix}u_{n+1}\\u_n\end{pmatrix}$.
    \item Déterminer la forme de Jordan de~$A$.
    \item Calculer $A^n$ et en déduire $u_n$.
  \end{enumerate}
\end{exercise}


% ============================================================
\section*{Résumé du chapitre}
\addcontentsline{toc}{section}{Résumé du chapitre}
% ============================================================

\begin{itemize}
  \item Un endomorphisme $f$ est \textbf{nilpotent}\index{nilpotent}
    si $f^\nu=0$ pour un certain $\nu\geq 1$. L'indice de nilpotence
    est le plus petit tel~$\nu$, et il vérifie $\nu\leq\dim E$.
    Le polynôme caractéristique est $X^n$ et le polynôme minimal
    est $X^\nu$.

  \item Un \textbf{sous-espace caractéristique}\index{sous-espace
    caractéristique} (ou espace propre généralisé) est
    $F_\lambda(f)=\Ker((f-\lambda\Id)^n)$. Sa dimension égale la
    multiplicité algébrique de~$\lambda$.

  \item L'espace $E$ se décompose en somme directe de sous-espaces
    caractéristiques~:
    $E=\bigoplus_{\lambda\in\Spec(f)}F_\lambda(f)$.

  \item Un \textbf{bloc de Jordan}\index{bloc de Jordan} $J_k(\lambda)$
    est une matrice $k\times k$ avec $\lambda$ sur la diagonale et des
    $1$ sur la sur-diagonale.

  \item \textbf{Théorème de Jordan~:}\index{forme normale de Jordan}
    toute matrice sur~$\C$ est semblable à une matrice de Jordan,
    unique à l'ordre des blocs près.

  \item Le \textbf{polynôme minimal}\index{polynôme minimal!et Jordan}
    est $\mu_f(X)=\prod_i(X-\lambda_i)^{\nu_i}$ où $\nu_i$ est la
    taille du plus grand bloc de Jordan associé à~$\lambda_i$.

  \item Une matrice est \textbf{diagonalisable} si et seulement si tous
    ses blocs de Jordan sont de taille~$1$.

  \item Le \textbf{calcul des tailles des blocs}\index{blocs de Jordan!calcul}
    se fait via les dimensions $\dim\Ker((A-\lambda I)^k)$
    pour $k=1,2,\ldots$

  \item L'\textbf{exponentielle de matrice}\index{exponentielle de matrice}
    $e^{tA}$ se calcule bloc par bloc via la forme de Jordan, ce qui
    permet de résoudre le système $X'=AX$.

  \item La \textbf{décomposition de Dunford}\index{décomposition de Dunford}
    $f=d+n$ (diagonalisable + nilpotent, commutants) est une conséquence
    directe de la forme de Jordan.
\end{itemize}
